\PassOptionsToPackage{unicode=true}{hyperref} % options for packages loaded elsewhere
\PassOptionsToPackage{hyphens}{url}
%
\documentclass[10pt,]{article}
\usepackage{lmodern}
\usepackage{amssymb,amsmath}
\usepackage{ifxetex,ifluatex}
\usepackage{fixltx2e} % provides \textsubscript
\ifnum 0\ifxetex 1\fi\ifluatex 1\fi=0 % if pdftex
  \usepackage[T1]{fontenc}
  \usepackage[utf8]{inputenc}
  \usepackage{textcomp} % provides euro and other symbols
\else % if luatex or xelatex
  \usepackage{unicode-math}
  \defaultfontfeatures{Ligatures=TeX,Scale=MatchLowercase}
\fi
% use upquote if available, for straight quotes in verbatim environments
\IfFileExists{upquote.sty}{\usepackage{upquote}}{}
% use microtype if available
\IfFileExists{microtype.sty}{%
\usepackage[]{microtype}
\UseMicrotypeSet[protrusion]{basicmath} % disable protrusion for tt fonts
}{}
\IfFileExists{parskip.sty}{%
\usepackage{parskip}
}{% else
\setlength{\parindent}{0pt}
\setlength{\parskip}{6pt plus 2pt minus 1pt}
}
\usepackage{hyperref}
\hypersetup{
            pdftitle={Holmes (2005), supplementary equations (8)--(17): closed-form solutions, relation to the gravestone Volterra equation, and misprints},
            pdfborder={0 0 0},
            breaklinks=true}
\urlstyle{same}  % don't use monospace font for urls
\usepackage[margin=2.2cm]{geometry}
\usepackage{longtable,booktabs}
% Fix footnotes in tables (requires footnote package)
\IfFileExists{footnote.sty}{\usepackage{footnote}\makesavenoteenv{longtable}}{}
\setlength{\emergencystretch}{3em}  % prevent overfull lines
\providecommand{\tightlist}{%
  \setlength{\itemsep}{0pt}\setlength{\parskip}{0pt}}
\setcounter{secnumdepth}{0}
% Redefines (sub)paragraphs to behave more like sections
\ifx\paragraph\undefined\else
\let\oldparagraph\paragraph
\renewcommand{\paragraph}[1]{\oldparagraph{#1}\mbox{}}
\fi
\ifx\subparagraph\undefined\else
\let\oldsubparagraph\subparagraph
\renewcommand{\subparagraph}[1]{\oldsubparagraph{#1}\mbox{}}
\fi

% set default figure placement to htbp
\makeatletter
\def\fps@figure{htbp}
\makeatother

\usepackage{amssymb}
\usepackage{newunicodechar}
\newunicodechar{′}{\ensuremath{'}}
\newunicodechar{→}{\ensuremath{\to}}
\newunicodechar{∎}{\ensuremath{\square}}
\allowdisplaybreaks

\title{Holmes (2005), supplementary equations (8)--(17): closed-form solutions,
relation to the gravestone Volterra equation, and misprints}
\date{}

\begin{document}
\maketitle

This note merges three independent attacks on the problem: route A
(generating function), route B (score identities) and route C
(numerics). Every identity used below is checked symbolically by
\texttt{verify\_all.py}; its output is in
\texttt{verify\_all\_output.txt} (172 checks, all reduce to exactly
\texttt{0} or \texttt{True}). The numerical evidence for Q3 is in
\texttt{numeric\_check.py} and \texttt{numeric\_check\_output.txt};
those numbers are checks, not part of any proof.

Sources:

\begin{itemize}
\tightlist
\item
  {[}H-S{]} the supplement of Holmes (2005), \emph{Bioinformatics}
  21:2294, pp.~9--14, read from the page images.
\item
  {[}H-M{]} the main text, PMC version, §3.3.
\item
  {[}DART{]} \texttt{src/tkf/tkfcoeff.cc} in github.com/ihh/dart,
  Holmes's implementation. The lines used here are unchanged since the
  earliest CVS revision (2003-11-10; see
  \texttt{dart\_tkfcoeff\_note.md}). What it computes is summarised in
  §5.3.
\item
  {[}P{]} paper 1 (the MBE revision); \texttt{eq:} names are its LaTeX
  equation labels.
\end{itemize}

\hypertarget{answers}{%
\section{0. Answers}\label{answers}}

\textbf{Q1 (closed forms).} For each zone type
\(z\in\{\circ,\bullet,\circ\bullet,\bullet\bullet\}\), the closed forms
\(I_z\) (insertions) and \(Z_z\) (exposure) are defined from the TKF91
pair-HMM in two equivalent ways:

\begin{itemize}
\tightlist
\item
  by the score identities eq:score-bds plus the conservation law
  eq:conservation, applied to a factorisation of the path weight into
  zone weights (Definition 2);
\item
  as logarithmic derivatives, in the gravestone variable \(w\) or in an
  exposure tilt \(\theta\), of three closed-form functions \(B,W,R\)
  obtained from eq:pgf-closed and eq:ID-closed (§6, Proposition 6).
\end{itemize}

Three theorems settle Q1.

\begin{itemize}
\tightlist
\item
  \textbf{Theorem 1.} For every alignment path,
  \(\sum_zN_zI_z=\mathbb E[B]\), \(\sum_zN_z(I_z-c_z)=\mathbb E[D]\) and
  \(\sum_zN_zZ_z=\mathbb E[S]+T\). No other per-zone coefficients have
  this property.
\item
  \textbf{Theorem 2.} The closed forms satisfy (8), (9)--(16) as
  printed, and (17) once the two ``\(t+\)'' terms are deleted. This
  holds for all \(\lambda,\mu>0\) including \(\lambda=\mu\).
\item
  \textbf{Theorem 3.} Each of these equations has exactly one bounded
  solution. So the exact solution of (8)--(16), (17′) is the closed
  forms, and an E-step built on it is exact. Holmes's first-order scheme
  (the {[}DART{]} port) is observed numerically to converge to the
  closed forms, at first order.
\end{itemize}

\textbf{Q2 (relation to eq:volterra).} The hypothesis holds, with two
refinements.

\begin{itemize}
\tightlist
\item
  \textbf{(9) and (14)} are first variations of the exposure-tilted
  eq:volterra at \(z=0\), in \(w\) and \(\theta\) respectively, divided
  by the fate probability. They are therefore linear Volterra equations
  of the second kind.
\item
  \textbf{(11) and (16)} are the \(w\)- and \(\theta\)-variations of the
  \(z^1\) coefficient of the death term, after division by the factor
  \(W\).
\item
  \textbf{(10), (15), (12) and (17′)} are not integral equations. They
  are variations of \(\mathcal I_D=\exp(\lambda\int(f-1))\) and of its
  geometric ratio, so they reduce to quadratures.
\item
  \textbf{Propagation event versus \(D\to L\,D\).} The production
  \(D\to L\cdot D\) peels off the most recent birth of any kind.
  Holmes's propagation event is the most recent birth whose family has a
  survivor. Summing \(D\to L\cdot D\) over the extinct children born
  after it gives \(\mathcal I_D(T-t,T-t;0)=1-\beta(T-t)\), the numerator
  of \(c\). The same form over the earlier window \([0,t]\) is
  \(v(t\mid T)/e^{-\mu t}\) in (8). The two decompositions coincide
  after this regrouping.
\item
  \textbf{(9) and (14) are not last-event decompositions.} They come
  from differentiating the exponential \(\mathcal I_D\) (a first-moment
  formula), followed by an exchange of the order of integration between
  the grammar's death time \(\ell\) and the birth time \(t\).
\end{itemize}

\textbf{Q3 (misprints).} The supplement contains three errors. All other
displayed equations are correct: (8)--(16), (13), and the sums
\(\hat S,\hat I,\hat D\) of §5.4.

\begin{enumerate}
\def\labelenumi{\arabic{enumi}.}
\tightlist
\item
  \textbf{Eq. (17).} Both ``\(t+\)'' terms must be deleted. As printed,
  the right-hand side equals the correct \(Z_{\bullet\bullet}\) plus
  \(\mathbb E[t_{\mathcal P}\mid\bullet\bullet]=\beta(T)^{-1}\int_0^T\beta=(\lambda T+\log(1-\beta(T)))/(\mu\beta(T))>0\).
  (16), which has no \(t\), is the consistent one. Holmes's code
  {[}DART{]} computes \texttt{Z11} without the \(t\), and his own Fig. 9
  (bottom) allocates exposure as (16) and (17′) do (§5.1, evidence 5).
\item
  \textbf{§5.4, last line.} It should read
  \(\hat Z=\sum_zN_z\hat Z_z-T\). {[}H-M{]} §3.3 defines \(\hat Z\) as
  the time accrued by \emph{mortal} links, but the zone sum also
  contains the immortal link's \(T\), through (15) applied to the
  S-zone. {[}DART{]} subtracts \(T\) (``subtract site-time contribution
  of immortal link'') and sets \(\lambda=I/(Z+T)\), \(\mu=D/Z\), which
  is eq:mstep.
\item
  \textbf{Fig. 9 caption, bottom example.} ``Ancestor survives'' should
  read ``Ancestor dies''. The drawing shows the ancestor deleted, and
  the first zone is \(\circ\bullet\). (\(q_0\), used before (11), is not
  an error: the main article's Eq. (2) defines \(q_n(t)\), and
  \(q_0=(1-\alpha)(1-\gamma)=\mu\beta/\lambda\), the extinction
  probability. The same factor appears in (13) and in Fig. 8.)
\end{enumerate}

\hypertarget{setting-and-assumptions}{%
\section{1. Setting and assumptions}\label{setting-and-assumptions}}

TKF91 has insertion rate \(\lambda\) and deletion rate \(\mu\), with
\(\lambda,\mu>0\).

\begin{itemize}
\tightlist
\item
  \textbf{Branch length and links.} \(T\) is the branch length (the
  paper's \(t\)). Each link alive at time \(\tau\) inserts a child
  immediately to its right at rate \(\lambda\). Mortal links die at rate
  \(\mu\); the immortal link never dies.
\item
  \textbf{Factors.} \(\alpha(s)=e^{-\mu s}\),
  \(\eta(s)=e^{-\lambda s}\), \(\delta=\mu\eta-\lambda\alpha\),
  \(\beta=\lambda(\eta-\alpha)/\delta\),
  \(\gamma=1-\mu\beta/(\lambda(1-\alpha))\), as in eq:abg. These are
  Holmes's \(\alpha,\beta,\gamma\). Unless an argument is shown,
  \(\alpha=\alpha(T)\) and so on.
\item
  \textbf{Holmes's rates} (A1 below supplies \(q_0\)):
  \[\lambda_v=\lambda\alpha,\qquad\lambda_w=\lambda(1-\alpha)\gamma,\qquad\lambda_x=\lambda(1-\alpha)(1-\gamma)=\lambda q_0 .\]
\item
  \textbf{Abbreviations.} \(k=\lambda-\mu\) and
  \(\epsilon=\lambda(1-\alpha)-\mu(1-\eta)\). Inside integrals,
  \(s=T-t\).
\item
  \textbf{Notation from {[}P{]}.} \(\mathcal I_D(t,T;z,w)\) keeps the
  paper's argument order: remaining time, then window length. \(T\) is
  also the branch length, so \(\mathcal I_D(T,t;0)\) is the window
  \([0,t]\) of a link with \(T\) remaining. Unlike eq:pgf-closed, the
  note does not require \(z_+\) to be the larger root; the closed forms
  are invariant under \(z_+\leftrightarrow z_-\) with
  \(\Delta\to-\Delta\) (§6.1).
\end{itemize}

\textbf{The equations as printed} ({[}H-S{]} pp.~9--11, transcribed from
the page images; hats as printed). Holmes defines
\(\lambda_v(t)=\lambda\alpha(t)\),
\(\lambda_w(t)=\lambda(1-\alpha(t))\gamma(t)\),
\(\lambda_x(t)=\lambda(1-\alpha(t))(1-\gamma(t))\), and: \begin{align*}
\log v(t\mid T)&=-\int_0^t\big(\mu+\lambda_v(T-t')+\lambda_w(T-t')\big)\,dt' \tag*{(8)}\\
\hat I_\circ(T)&=\int_{t=0}^T a(t\mid T)\,\lambda_x(T-t)\big(1+\hat I_\circ(T-t)\big)\,dt \tag*{(9)}\\
\hat I_\bullet(T)&=\int_{t=0}^T \lambda_x(T-t)\big(1+\hat I_\circ(T-t)\big)\,dt \tag*{(10)}\\
\hat I_{\circ\bullet}(T)&=\int_{t=0}^T b(t\mid T)\Big(\lambda_v(T-t)\big(1+\hat I_\circ(T-t)\big)\\
&\qquad\qquad+\lambda_w(T-t)\big(1+\hat I_\circ(T-t)+\hat I_{\circ\bullet}(T-t)\big)\Big)\,dt \tag*{(11)}\\
\hat I_{\bullet\bullet}(T)&=\int_{t=0}^T c(t\mid T)\Big(\lambda_v(T-t)\big(1+\hat I_\bullet(T-t)\big)\\
&\qquad\qquad+\lambda_w(T-t)\big(1+\hat I_\bullet(T-t)+\hat I_{\circ\bullet}(T-t)\big)\Big)\,dt \tag*{(12)}\\
a(t\mid T)&=\frac{v(t\mid T)\,(1-\alpha(T-t))(1-\gamma(T-t))}{(1-\alpha(T))(1-\gamma(T))},\qquad
b(t\mid T)=\frac{\alpha(t)\,(1-\alpha(T-t))(1-\gamma(T-t))}{(1-\alpha(T))\gamma(T)},\\
c(t\mid T)&=\frac{1-\beta(T-t)}{\beta(T)} \tag*{(13)}\\
\hat Z_\circ(T)&=\int_{t=0}^T a(t\mid T)\big(1+\lambda_x(T-t)\hat Z_\circ(T-t)\big)\,dt \tag*{(14)}\\
\hat Z_\bullet(T)&=\int_{t=0}^T \big(1+\lambda_x(T-t)\hat Z_\circ(T-t)\big)\,dt \tag*{(15)}\\
\hat Z_{\circ\bullet}(T)&=\int_{t=0}^T b(t\mid T)\Big(\lambda_v(T-t)\hat Z_\circ(T-t)\\
&\qquad\qquad+\lambda_w(T-t)\big(\hat Z_\circ(T-t)+\hat Z_{\circ\bullet}(T-t)\big)\Big)\,dt \tag*{(16)}\\
\hat Z_{\bullet\bullet}(T)&=\int_{t=0}^T c(t\mid T)\Big(\lambda_v(T-t)\big(t+\hat Z_\bullet(T-t)\big)\\
&\qquad\qquad+\lambda_w(T-t)\big(t+\hat Z_\bullet(T-t)+\hat Z_{\circ\bullet}(T-t)\big)\Big)\,dt \tag*{(17)}
\end{align*} (17′) denotes (17) with both ``\(t+\)'' deleted. The note's
closed forms \(I_z,Z_z\) carry no hats.

\textbf{A1 (\(q_0\)).} \(q_0(s)=(1-\alpha(s))(1-\gamma(s))\) ({[}H-M{]}
Eq. (2), \(q_n\) at \(n=0\)), the probability that a link with \(s\)
time remaining dies and leaves no surviving descendant. This is the
factor in (13) and in Fig. 8(i),(iii). We write \(q\) for \(q_0\).

\textbf{A2 (zones).} Zone types are those of Holmes's transition lists:

\begin{itemize}
\tightlist
\item
  \(\circ\): DM, DD, DE;
\item
  \(\bullet\): SM, SD, SE, MM, MD, ME, IM, ID, IE;
\item
  \(\circ\bullet\): DI;
\item
  \(\bullet\bullet\): SI, MI, II.
\end{itemize}

So a zone runs from its marker to the next \emph{observed} link, which
may be a deleted ancestral residue (route C, note 5).

\textbf{A3 (exposure).} \(S\) is the exposure of mortal links (the
paper's \(S\)). ``\(Z\)'' is used only for Holmes's quantities.

\textbf{Lemma 0} (verify\_all §2--3).

\begin{enumerate}
\def\labelenumi{(\alph{enumi})}
\item
  \(\lambda_x=\mu\beta\), i.e.~\(q=\mu\beta/\lambda\).
\item
  \(\lambda_v+\lambda_w+\lambda_x=\lambda\).
\item
  \(\beta'=(1-\beta)(\lambda-\mu\beta)\), and hence
  \((1-\beta)(\lambda_v+\lambda_w)=\beta'\).
\item
  \(q'=\mu(1-\beta)(\lambda-\mu\beta)/\lambda\ge0\), because
  \(0\le\beta<\min(1,\lambda/\mu)\).
\item
  As \(T\to0\), \((1-\alpha)\gamma\sim\lambda\mu T^2/2\) and
  \(q\sim\mu T\).
\end{enumerate}

\hypertarget{what-must-equal-what-q1-first-half}{%
\section{2. What must equal what (Q1, first
half)}\label{what-must-equal-what-q1-first-half}}

\textbf{Lemma 1 (zone bookkeeping).} Let \(N_z\) be the zone counts of a
path \(\mathrm S\to\dots\to\mathrm E\). Then
\[\#I=N_{\circ\bullet}+N_{\bullet\bullet},\quad \#D=N_\circ+N_{\circ\bullet},\quad \#M=N_\bullet-N_{\circ\bullet}-1,\quad i=N_\circ+N_\bullet-1,\quad j-i=N_{\bullet\bullet}-N_\circ .\]
The weight of the path in eq:cond\_tkf91 is
\[P(\pi)=\alpha^{-1}\prod_zP_z^{N_z},\qquad P_\circ=(1-\alpha)(1-\gamma),\ \ P_\bullet=\alpha(1-\beta),\ \ P_{\circ\bullet}=\frac{(1-\alpha)\gamma}{\alpha},\ \ P_{\bullet\bullet}=\beta .\]

\emph{Proof.}

\begin{itemize}
\tightlist
\item
  Every transition out of D is either DI (\(\circ\bullet\)) or one of
  DM, DD, DE (\(\circ\)). Every transition into I is \(\circ\bullet\) or
  \(\bullet\bullet\).
\item
  Transitions out of S, M or I number
  \(N_\bullet+N_{\bullet\bullet}=1+\#M+\#I\).
\item
  Each weight in eq:cond\_tkf91 is a source-row factor (\(\beta\) or
  \(1-\beta\) from S, M, I; \(\gamma\) or \(1-\gamma\) from D) times a
  destination factor (\(\alpha\) for M, \(1-\alpha\) for D, \(1\) for I
  or E). Collect powers.
\end{itemize}

Both statements are checked exhaustively on all 3280 paths of length
\(\le7\) (verify\_all §1). ∎

The factor \(\alpha^{-1}=e^{\mu T}\) and the \(\alpha^{-1}\) inside
\(P_{\circ\bullet}\) are the bookkeeping behind Holmes's ``WLOG the
marker is ancestral''. Every \(\bullet\) zone carries one \(\alpha\), as
if its marker were a surviving ancestral link. The path has only
\(\#M=N_\bullet-N_{\circ\bullet}-1\) such factors. The deficit is one
per \(\circ\bullet\) zone and one for the immortal link.

\textbf{Definition 2 (our closed forms).} Put
\[\ell_z=\log P_z,\quad L_z=\lambda\partial_\lambda\ell_z,\quad M_z=\mu\partial_\mu\ell_z\ (T\text{ fixed}),\quad c_\circ=-1,\ c_\bullet=c_{\circ\bullet}=0,\ c_{\bullet\bullet}=+1,\]
\[\boxed{\;Z_z=\frac{c_z+M_z-L_z}{\lambda-\mu},\qquad I_z=L_z+\lambda Z_z,\qquad D_z=I_z-c_z=M_z+\mu Z_z\;}\]
These are eq:score-bds and eq:conservation applied zone by zone: the
per-zone analogue of eq:ES without its immortal-link term
\(-\lambda t\). By Theorem 1 the zone sum \(\sum_zN_zZ_z\) equals eq:ES
plus \(T\).

\textbf{Theorem 1 (exact totals).} For every path \(\pi\),
\[\sum_zN_zI_z=\mathbb E[B\mid\pi],\qquad\sum_zN_zD_z=\mathbb E[D\mid\pi],\qquad\sum_zN_zZ_z=\mathbb E[S\mid\pi]+T,\]
and so also for Baum--Welch expected counts, by linearity.

\emph{Proof.}

\begin{itemize}
\tightlist
\item
  The score identity holds for any event, in particular for \(\pi\).
  With Lemma 1, \(\lambda\partial_\lambda\log P(\pi)=\sum N_zL_z\),
  because \(\lambda\partial_\lambda\log\alpha^{-1}=0\). Also
  \(\mu\partial_\mu\log P(\pi)=\sum N_zM_z+\mu T\).
\item
  eq:score-bds gives
  \(\sum N_zL_z=\mathbb E[B]-\lambda(\mathbb E[S]+T)\) and
  \(\sum N_zM_z+\mu T=\mathbb E[D]-\mu\mathbb E[S]\). Lemma 1 and
  eq:conservation give
  \(\mathbb E[B]-\mathbb E[D]=N_{\bullet\bullet}-N_\circ=\sum N_zc_z\).
\item
  Solving these three linear equations gives the claim. ∎
\end{itemize}

(verify\_all §2, ``Theorem 1''.)

\textbf{Equal rates.} Definition 2 divides by \(\lambda-\mu\). For
\(\lambda=\mu\), \(I_z\) and \(Z_z\) are defined as the \(k\to0\)
limits, which have no pole (verify\_all §4; values in §3). Theorem 1
then holds at \(\lambda=\mu\) by continuity: each
\(\mathbb E[\cdot\mid\pi]\) is a ratio of functions analytic in
\(\lambda\) whose denominator \(P(\pi)\) is positive.

\textbf{Uniqueness of the coefficients.} The vectors
\((1,N_\circ,N_\bullet,N_{\circ\bullet},N_{\bullet\bullet})\) for the
paths \(\varnothing\), M, D, I, DI are linearly independent; their
determinant is nonzero, and the rank over all paths of length \(\le7\)
is 5. Hence an affine function \(K+\sum_zN_zk_z\) of the zone counts
determines \(K\) and every \(k_z\). Holmes's sums
\(\hat I=\sum N_z\hat I_z\) and \(\hat Z=\sum N_z\hat Z_z\) have
\(K=0\), which gives two consequences:

\begin{itemize}
\tightlist
\item
  \(\hat I=\mathbb E[B]\) for all data \textbf{iff} \(\hat I_z=I_z\) for
  all \(z\); then
  \(\hat D=\hat I+N_\circ-N_{\bullet\bullet}=\mathbb E[D]\).
\item
  \(\sum N_z\hat Z_z=\mathbb E[S]+T\) for all data \textbf{iff}
  \(\hat Z_z=Z_z\). No constant-free sum can equal \(\mathbb E[S]\),
  whose representation has constant \(K=-T\). Since {[}H-M{]} defines
  \(\hat Z\) as mortal-link time, §5.4 needs the ``\(-T\)'' (misprint
  2).
\end{itemize}

So the per-zone targets are pinned down. The constant
\(\alpha^{-1}=e^{\mu T}\) in Lemma 1 removes the death hazard that the
immortal link's \(\bullet\) zone would otherwise carry. The zone sum
therefore counts the immortal link's exposure \(T\). That exposure pairs
with the paper's \(-\lambda t\) (the \(\lambda\)-exposure is \(S+t\))
but not with the \(\mu\) side (the \(\mu\)-exposure is \(S\)). Hence
\(\sum_zN_zZ_z=\mathbb E[S]+T\), and eq:mstep divides \(B\) by \(S+T\)
and \(D\) by \(S\).

\textbf{Family form (route A, same numbers).} Group the zones by the
ancestral link (or the immortal link) and its run of \(n\) surviving
descendants:

\begin{longtable}[]{@{}ll@{}}
\toprule
\begin{minipage}[b]{0.47\columnwidth}\raggedright
family\strut
\end{minipage} & \begin{minipage}[b]{0.47\columnwidth}\raggedright
zones\strut
\end{minipage}\tabularnewline
\midrule
\endhead
\begin{minipage}[t]{0.47\columnwidth}\raggedright
survives, \(n\)\strut
\end{minipage} & \begin{minipage}[t]{0.47\columnwidth}\raggedright
\(n\times\bullet\bullet\), then one \(\bullet\)\strut
\end{minipage}\tabularnewline
\begin{minipage}[t]{0.47\columnwidth}\raggedright
dies, \(n=0\)\strut
\end{minipage} & \begin{minipage}[t]{0.47\columnwidth}\raggedright
one \(\circ\)\strut
\end{minipage}\tabularnewline
\begin{minipage}[t]{0.47\columnwidth}\raggedright
dies, \(n\ge1\)\strut
\end{minipage} & \begin{minipage}[t]{0.47\columnwidth}\raggedright
one \(\circ\bullet\), then \((n-1)\times\bullet\bullet\), then one
\(\bullet\)\strut
\end{minipage}\tabularnewline
\bottomrule
\end{longtable}

The family probabilities are \(\alpha(1-\beta)\beta^n\), \(q\) and
\((1-\alpha)\gamma(1-\beta)\beta^{n-1}\). Applying the score identities
to these events gives the following (verify\_all §5, ``family'', for
\(n\le2\)):

\begin{itemize}
\tightlist
\item
  \(X_\bullet=\mathbb E[X_{\rm fam}\mid\text{surv},0]\)
\item
  \(X_\circ=\mathbb E[X_{\rm fam}\mid\text{die},0]\)
\item
  \(X_{\bullet\bullet}=\mathbb E[X_{\rm fam}\mid\text{surv},n{+}1]-\mathbb E[X_{\rm fam}\mid\text{surv},n]\);
  the ``die'' families give the same increment.
\item
  \(X_{\circ\bullet}=\mathbb E[X_{\rm fam}\mid\text{die},1]-X_\bullet\)
\end{itemize}

Here \(X\in\{B,S\}\) and \(S_{\rm fam}\) includes the founder's own
exposure. The immortal family behaves like a surviving mortal family
with exposure \(T\) for the founder, which is again the ``\(+T\)''.

\hypertarget{the-closed-forms}{%
\section{3. The closed forms}\label{the-closed-forms}}

Factored forms (verify\_all §2; each is checked equal to Definition 2,
and route A's and route C's forms are checked equal to these): \[
\begin{aligned}
I_\circ&=\frac{\lambda\mu\,[\eta^2-\alpha^2+2kT\eta\alpha]}{k(\alpha-\eta)\delta}, &
Z_\circ&=\frac{\mu\eta^2-\lambda\alpha^2+k\eta\alpha+k(\lambda+\mu)T\eta\alpha}{k(\alpha-\eta)\delta},\\
I_\bullet&=\frac{\lambda\mu\,[(\lambda+\mu)(\alpha-\eta)-kT(\lambda\alpha+\mu\eta)]}{k^2\delta}, &
Z_\bullet&=\frac{2\lambda\mu(\alpha-\eta)-kT(\lambda^2\alpha+\mu^2\eta)}{k^2\delta},\\
I_{\circ\bullet}&=\frac{k}{\epsilon}\,I_\bullet, &
Z_{\circ\bullet}&=\frac{\lambda\mu\,[2(\alpha-\eta)-kT(\alpha+\eta)]}{k\,\delta\,\epsilon},\\
I_{\bullet\bullet}&=1+I_\circ, & Z_{\bullet\bullet}&=Z_\circ .
\end{aligned}
\] The last line follows from
\(P_\circ=(1-\alpha)(1-\gamma)=(\mu/\lambda)\beta=(\mu/\lambda)P_{\bullet\bullet}\)
(Lemma 0a). The constant factor \(\mu/\lambda\) has
\(\lambda\partial_\lambda\log(\mu/\lambda)=-1\) and
\(\mu\partial_\mu\log(\mu/\lambda)=+1\), so
\(L_\circ=L_{\bullet\bullet}-1\) and \(M_\circ=M_{\bullet\bullet}+1\).
In Definition 2 this raises \(M_z-L_z\) by exactly
\(2=c_{\bullet\bullet}-c_\circ\), so \(Z_\circ=Z_{\bullet\bullet}\), and
then \(I_\circ=L_\circ+\lambda Z_\circ=I_{\bullet\bullet}-1\). (§6.2
states the same fact for the generating functions.)

\textbf{Equal rates} (\(\sigma=\mu T\); the \(\lambda\to\mu\) limits
have no pole in \(k\), verify\_all §4):

\begin{longtable}[]{@{}ll@{}}
\toprule
\begin{minipage}[b]{0.47\columnwidth}\raggedright
quantity\strut
\end{minipage} & \begin{minipage}[b]{0.47\columnwidth}\raggedright
value at \(\lambda=\mu\)\strut
\end{minipage}\tabularnewline
\midrule
\endhead
\begin{minipage}[t]{0.47\columnwidth}\raggedright
\(I_\circ\)\strut
\end{minipage} & \begin{minipage}[t]{0.47\columnwidth}\raggedright
\(\dfrac{\sigma^2}{3(1+\sigma)}\)\strut
\end{minipage}\tabularnewline
\begin{minipage}[t]{0.47\columnwidth}\raggedright
\(Z_\circ=Z_{\bullet\bullet}\)\strut
\end{minipage} & \begin{minipage}[t]{0.47\columnwidth}\raggedright
\(\dfrac{T(2\sigma+3)}{6(1+\sigma)}\)\strut
\end{minipage}\tabularnewline
\begin{minipage}[t]{0.47\columnwidth}\raggedright
\(I_\bullet\)\strut
\end{minipage} & \begin{minipage}[t]{0.47\columnwidth}\raggedright
\(\dfrac{\sigma^2(\sigma+3)}{6(1+\sigma)}\)\strut
\end{minipage}\tabularnewline
\begin{minipage}[t]{0.47\columnwidth}\raggedright
\(Z_\bullet\)\strut
\end{minipage} & \begin{minipage}[t]{0.47\columnwidth}\raggedright
\(\dfrac{T(\sigma^2+6\sigma+6)}{6(1+\sigma)}\)\strut
\end{minipage}\tabularnewline
\begin{minipage}[t]{0.47\columnwidth}\raggedright
\(I_{\circ\bullet}\)\strut
\end{minipage} & \begin{minipage}[t]{0.47\columnwidth}\raggedright
\(\dfrac{\sigma^2(\sigma+3)}{6(1+\sigma)(1-(1+\sigma)\alpha)}\)\strut
\end{minipage}\tabularnewline
\begin{minipage}[t]{0.47\columnwidth}\raggedright
\(Z_{\circ\bullet}\)\strut
\end{minipage} & \begin{minipage}[t]{0.47\columnwidth}\raggedright
\(\dfrac{\sigma^2T}{6(1+\sigma)(1-(1+\sigma)\alpha)}\)\strut
\end{minipage}\tabularnewline
\begin{minipage}[t]{0.47\columnwidth}\raggedright
\(I_{\bullet\bullet}\)\strut
\end{minipage} & \begin{minipage}[t]{0.47\columnwidth}\raggedright
\(1+I_\circ\)\strut
\end{minipage}\tabularnewline
\bottomrule
\end{longtable}

\textbf{Floating point near \(\lambda=\mu\).} The factored forms divide
by powers of \(k\) and cancel catastrophically in float64. At \(T=1\)
the relative error is about \(10^{-6}\) at \(|\lambda-\mu|=10^{-3}\) and
about \(10^{-3}\) at \(|\lambda-\mu|=10^{-4}\). When
\(|\lambda-\mu|T\lesssim10^{-3}\), use a series in \(k\) (the equal-rate
table is its leading term) or forms rewritten with \texttt{expm1}.

\textbf{Posterior objects.} Lemma 0(c) gives
\[v(t\mid T)=\frac{\alpha(t)\,(1-\beta(T))}{1-\beta(s)}=e^{-\mu t}\,\mathcal I_D(T,t;0,1),\]
the second form being eq:ID-closed at \(z=0\), \(w=1\). With \(F\) and
\(G\) denoting separable factors,
\[a(t\mid T)=\frac{v(t\mid T)q(s)}{q(T)}=F_a(T)G_a(s),\qquad b(t\mid T)=\frac{\alpha(t)q(s)}{(1-\alpha)\gamma}=F_b(T)G_b(s),\qquad c(t\mid T)=\frac{1-\beta(s)}{\beta(T)}=F_c(T)G_c(s),\]
\[F_a=\frac{\alpha(1-\beta)}{q},\ G_a=\frac1{F_a};\qquad F_b=\frac{\alpha}{(1-\alpha)\gamma},\ G_b=\frac{q}{\alpha};\qquad F_c=\frac1\beta,\ G_c=1-\beta .\]

\textbf{Remark (symmetry).} All eight \(I_z,Z_z\) are invariant under
\(\lambda\leftrightarrow\mu\) (verify\_all §6). This is not used in any
proof.

\hypertarget{the-closed-forms-solve-holmess-equations-q1-second-half}{%
\section{4. The closed forms solve Holmes's equations (Q1, second
half)}\label{the-closed-forms-solve-holmess-equations-q1-second-half}}

\textbf{Lemma 4 (separable Volterra equations).} Let \(X\) be continuous
on \((0,T_{\max}]\), \(F>0\) and \(G R\) integrable near 0. Then
\(X(T)=\int_0^TF(T)G(T-t)R(T-t)\,dt\) for all \(T\) \textbf{iff}
\[\frac{d}{dT}\frac{X}{F}=G(T)R(T)\quad\text{and}\quad\lim_{T\to0^+}\frac XF=0 .\]
The proof is immediate after substituting \(s=T-t\).

\textbf{Theorem 2.} With \(v,a,b,c\) as in §3 and \(I_z,Z_z\) from
Definition 2:

\begin{longtable}[]{@{}ll@{}}
\toprule
\begin{minipage}[b]{0.47\columnwidth}\raggedright
eq.\strut
\end{minipage} & \begin{minipage}[b]{0.47\columnwidth}\raggedright
status\strut
\end{minipage}\tabularnewline
\midrule
\endhead
\begin{minipage}[t]{0.47\columnwidth}\raggedright
(8)\strut
\end{minipage} & \begin{minipage}[t]{0.47\columnwidth}\raggedright
holds:
\(\partial_t\log v=-\mu-\beta'(s)/(1-\beta(s))=-(\mu+\lambda_v+\lambda_w)(s)\)
by Lemma 0(b,c), and \(v(0\mid T)=1\)\strut
\end{minipage}\tabularnewline
\begin{minipage}[t]{0.47\columnwidth}\raggedright
(13)\strut
\end{minipage} & \begin{minipage}[t]{0.47\columnwidth}\raggedright
correct as printed (with A1). \(a\) is the posterior survival function
of the marker given \(\circ\): by (S) below,
\(\int_t^T\mu v(\ell\mid T)d\ell=v(t\mid T)q(s)\). Also
\(\int_0^Tb(\lambda_v+\lambda_w)=\int_0^Tc(\lambda_v+\lambda_w)=1\)\strut
\end{minipage}\tabularnewline
\begin{minipage}[t]{0.47\columnwidth}\raggedright
(9), (10), (11), (12)\strut
\end{minipage} & \begin{minipage}[t]{0.47\columnwidth}\raggedright
hold as printed\strut
\end{minipage}\tabularnewline
\begin{minipage}[t]{0.47\columnwidth}\raggedright
(14), (15), (16)\strut
\end{minipage} & \begin{minipage}[t]{0.47\columnwidth}\raggedright
hold as printed\strut
\end{minipage}\tabularnewline
\begin{minipage}[t]{0.47\columnwidth}\raggedright
(17′)\strut
\end{minipage} & \begin{minipage}[t]{0.47\columnwidth}\raggedright
(17) with both ``\(t+\)'' deleted holds\strut
\end{minipage}\tabularnewline
\bottomrule
\end{longtable}

\emph{Proof.} Each equation of the transcription in §1 is of the form in
Lemma 4, with the following data.

\begin{longtable}[]{@{}lllll@{}}
\toprule
\begin{minipage}[b]{0.17\columnwidth}\raggedright
eq.\strut
\end{minipage} & \begin{minipage}[b]{0.17\columnwidth}\raggedright
\(X\)\strut
\end{minipage} & \begin{minipage}[b]{0.17\columnwidth}\raggedright
\(F\)\strut
\end{minipage} & \begin{minipage}[b]{0.17\columnwidth}\raggedright
\(G\)\strut
\end{minipage} & \begin{minipage}[b]{0.17\columnwidth}\raggedright
\(R(s)\)\strut
\end{minipage}\tabularnewline
\midrule
\endhead
\begin{minipage}[t]{0.17\columnwidth}\raggedright
(9)\strut
\end{minipage} & \begin{minipage}[t]{0.17\columnwidth}\raggedright
\(I_\circ\)\strut
\end{minipage} & \begin{minipage}[t]{0.17\columnwidth}\raggedright
\(F_a\)\strut
\end{minipage} & \begin{minipage}[t]{0.17\columnwidth}\raggedright
\(G_a\)\strut
\end{minipage} & \begin{minipage}[t]{0.17\columnwidth}\raggedright
\(\lambda_x(1+I_\circ)\)\strut
\end{minipage}\tabularnewline
\begin{minipage}[t]{0.17\columnwidth}\raggedright
(10)\strut
\end{minipage} & \begin{minipage}[t]{0.17\columnwidth}\raggedright
\(I_\bullet\)\strut
\end{minipage} & \begin{minipage}[t]{0.17\columnwidth}\raggedright
1\strut
\end{minipage} & \begin{minipage}[t]{0.17\columnwidth}\raggedright
1\strut
\end{minipage} & \begin{minipage}[t]{0.17\columnwidth}\raggedright
\(\lambda_x(1+I_\circ)\)\strut
\end{minipage}\tabularnewline
\begin{minipage}[t]{0.17\columnwidth}\raggedright
(11)\strut
\end{minipage} & \begin{minipage}[t]{0.17\columnwidth}\raggedright
\(I_{\circ\bullet}\)\strut
\end{minipage} & \begin{minipage}[t]{0.17\columnwidth}\raggedright
\(F_b\)\strut
\end{minipage} & \begin{minipage}[t]{0.17\columnwidth}\raggedright
\(G_b\)\strut
\end{minipage} & \begin{minipage}[t]{0.17\columnwidth}\raggedright
\(\lambda_v(1+I_\circ)+\lambda_w(1+I_\circ+I_{\circ\bullet})\)\strut
\end{minipage}\tabularnewline
\begin{minipage}[t]{0.17\columnwidth}\raggedright
(12)\strut
\end{minipage} & \begin{minipage}[t]{0.17\columnwidth}\raggedright
\(I_{\bullet\bullet}\)\strut
\end{minipage} & \begin{minipage}[t]{0.17\columnwidth}\raggedright
\(F_c\)\strut
\end{minipage} & \begin{minipage}[t]{0.17\columnwidth}\raggedright
\(G_c\)\strut
\end{minipage} & \begin{minipage}[t]{0.17\columnwidth}\raggedright
\(\lambda_v(1+I_\bullet)+\lambda_w(1+I_\bullet+I_{\circ\bullet})\)\strut
\end{minipage}\tabularnewline
\begin{minipage}[t]{0.17\columnwidth}\raggedright
(14)\strut
\end{minipage} & \begin{minipage}[t]{0.17\columnwidth}\raggedright
\(Z_\circ\)\strut
\end{minipage} & \begin{minipage}[t]{0.17\columnwidth}\raggedright
\(F_a\)\strut
\end{minipage} & \begin{minipage}[t]{0.17\columnwidth}\raggedright
\(G_a\)\strut
\end{minipage} & \begin{minipage}[t]{0.17\columnwidth}\raggedright
\(1+\lambda_xZ_\circ\)\strut
\end{minipage}\tabularnewline
\begin{minipage}[t]{0.17\columnwidth}\raggedright
(15)\strut
\end{minipage} & \begin{minipage}[t]{0.17\columnwidth}\raggedright
\(Z_\bullet\)\strut
\end{minipage} & \begin{minipage}[t]{0.17\columnwidth}\raggedright
1\strut
\end{minipage} & \begin{minipage}[t]{0.17\columnwidth}\raggedright
1\strut
\end{minipage} & \begin{minipage}[t]{0.17\columnwidth}\raggedright
\(1+\lambda_xZ_\circ\)\strut
\end{minipage}\tabularnewline
\begin{minipage}[t]{0.17\columnwidth}\raggedright
(16)\strut
\end{minipage} & \begin{minipage}[t]{0.17\columnwidth}\raggedright
\(Z_{\circ\bullet}\)\strut
\end{minipage} & \begin{minipage}[t]{0.17\columnwidth}\raggedright
\(F_b\)\strut
\end{minipage} & \begin{minipage}[t]{0.17\columnwidth}\raggedright
\(G_b\)\strut
\end{minipage} & \begin{minipage}[t]{0.17\columnwidth}\raggedright
\(\lambda_vZ_\circ+\lambda_w(Z_\circ+Z_{\circ\bullet})\)\strut
\end{minipage}\tabularnewline
\begin{minipage}[t]{0.17\columnwidth}\raggedright
(17′)\strut
\end{minipage} & \begin{minipage}[t]{0.17\columnwidth}\raggedright
\(Z_{\bullet\bullet}\)\strut
\end{minipage} & \begin{minipage}[t]{0.17\columnwidth}\raggedright
\(F_c\)\strut
\end{minipage} & \begin{minipage}[t]{0.17\columnwidth}\raggedright
\(G_c\)\strut
\end{minipage} & \begin{minipage}[t]{0.17\columnwidth}\raggedright
\(\lambda_vZ_\bullet+\lambda_w(Z_\bullet+Z_{\circ\bullet})\)\strut
\end{minipage}\tabularnewline
\bottomrule
\end{longtable}

\begin{itemize}
\tightlist
\item
  \textbf{Derivative condition.} Write every function of \(T\) as a
  rational function of \((\lambda,\mu,T,\eta,\alpha)\), and apply
  \(d/dT=\partial_T-\lambda\eta\partial_\eta-\mu\alpha\partial_\alpha\).
  The difference \(\frac{d}{dT}(X/F)-GR\) cancels to the zero rational
  function. That makes it an identity after substituting
  \(\eta=e^{-\lambda T}\), \(\alpha=e^{-\mu T}\).
\item
  \textbf{Initial condition.} \(\lim_{T\to0^+}X/F=0\), by series
  expansion.
\item
  \textbf{The kernel statements in (13)} are \(d/dT\) identities of the
  same kind:

  \begin{itemize}
  \tightlist
  \item
    \(\frac{d}{dT}[1/F_b]=G_b(\lambda_v+\lambda_w)\) and
    \(\frac{d}{dT}[1/F_c]=G_c(\lambda_v+\lambda_w)\), with
    \(1/F_b,1/F_c\to0\);
  \item
    (S): \(\partial_t[v(t\mid T)q(s)]=-\mu v(t\mid T)\), with
    \(q(0)=0\). This is the evidence that \(a\) is the posterior
    survival function of the marker given \(\circ\).
  \end{itemize}
\item
  \textbf{Equal rates.} For \(\lambda=\mu\) the same identities are
  verified directly, with \(\beta=\mu T/(1+\mu T)\) and the equal-rate
  forms of §3; they are also the \(k\to0\) limits.
\end{itemize}

(verify\_all §3--4.) ∎

\textbf{Remark (small \(T\); denominators).} As \(T\to0\),
\(I_\circ,I_\bullet,Z_\circ,Z_\bullet,Z_{\circ\bullet},Z_{\bullet\bullet}\to0\)
and \(I_{\circ\bullet},I_{\bullet\bullet}\to1\) (a \(\circ\bullet\) or
\(\bullet\bullet\) zone contains at least one insertion). So every
\(GR\) in the table is bounded near 0. For \(T>0\) and \(\lambda\ne\mu\)
no denominator in §3 vanishes:

\begin{itemize}
\tightlist
\item
  \(k\ne0\) and \(\alpha\ne\eta\);
\item
  \(\delta=\mu e^{-\lambda T}-\lambda e^{-\mu T}\ne0\), because
  \(e^{(\mu-\lambda)T}=\lambda/\mu\) has no solution with \(T>0\) (if
  \(\lambda>\mu\) the left side is below 1 and the right side above 1,
  and conversely);
\item
  \(\epsilon=\lambda\mu\,(g(\mu)-g(\lambda))\ne0\), where
  \(g(x)=(1-e^{-xT})/x=\int_0^Te^{-xs}ds\) is strictly decreasing; hence
  also \((1-\alpha)\gamma=-\alpha\epsilon/\delta\ne0\).
\end{itemize}

(verify\_all §3: the eight limits and the identity for
\((1-\alpha)\gamma\).)

\textbf{Theorem 3 (uniqueness).} Fix \(T_{\max}>0\) and \(v,a,b,c\).
Once the lower-numbered quantities on its right-hand side are fixed,
each of (9), (14), (11), (16) has at most one solution \(X\) on
\((0,T_{\max}]\) for which \(X/F\) is integrable near 0, with \(F=F_a\)
for (9), (14) and \(F=F_b\) for (11), (16). Every bounded solution
qualifies, since \(1/F_a=q/(\alpha(1-\beta))\sim\mu T\) and
\(1/F_b=(1-\alpha)\gamma/\alpha\sim\lambda\mu T^2/2\) (Lemma 0e). (10),
(12), (15) and (17′) are explicit quadratures. Hence the system
(8)--(16), (17′) has exactly one solution in this class, the closed
forms.

\emph{Proof.} Let \(d\) be the difference of two solutions and
\(u=d/F\). Substitute \(s=T-t\) in the homogeneous equation for \(d\).

\begin{itemize}
\tightlist
\item
  \textbf{(9) and (14).} \(u(T)=\int_0^T\lambda_x(s)\,u(s)\,ds\),
  because \(G_aF_a=1\).
\item
  \textbf{(11) and (16).} Only the \(\lambda_w\) term contains the
  unknown, and \(G_bF_b\lambda_w=\lambda q\) (verify\_all §3). So
  \(u(T)=\int_0^T\lambda q(s)\,u(s)\,ds\).
\item
  In both cases \(u(T)=\int_0^Tk(s)u(s)\,ds\) with \(k\) continuous and
  \(0\le k\le\lambda\). Put \(U(T)=\int_0^T|u|\), finite by hypothesis.
  Then \(U(0)=0\) and \(U'=|u|\le\lambda U\), so \(e^{-\lambda T}U\) is
  nonincreasing and nonnegative, hence zero. So \(u\equiv0\) and
  \(d\equiv0\). (Equivalently: \(u(0^+)=0\) and \(u'=ku\).) ∎
\end{itemize}

\textbf{Consequence.} The exact solution of (8)--(16), (17′) is the
closed forms, so by Theorem 1 an E-step built on it is exact:
\(\hat I=\mathbb E[B]\), \(\hat D=\mathbb E[D]\), and
\(\sum N_z\hat Z_z-T=\mathbb E[S]\). Uniqueness says nothing about a
discretisation. Holmes's own first-order scheme, ported from {[}DART{]}
in \texttt{numeric\_check.py}, is observed to converge to the closed
forms. At \((\lambda,\mu,T)=(1,0.5,1)\) and \((2,0.3,1)\), every
coefficient converges at first order: the error halves as \(N\) doubles,
and is at most \(5\times10^{-4}\) at \(N=2000\). \texttt{Z11} tends to
\(0.4224\) and \(0.3571\) respectively. The printed-(17) values are
\(1.018\) and \(1.013\).

\hypertarget{q3-misprints-and-why-16-rather-than-17}{%
\section{5. Q3: misprints, and why (16) rather than
(17)}\label{q3-misprints-and-why-16-rather-than-17}}

\hypertarget{eq.-17}{%
\subsection{5.1 Eq. (17)}\label{eq.-17}}

By Lemma 0(c),
\(c(t\mid T)(\lambda_v+\lambda_w)(s)=\beta'(s)/\beta(T)\). This is the
posterior density of the propagation time given a \(\bullet\bullet\)
zone. Therefore
\[\int_0^Tc(t\mid T)(\lambda_v+\lambda_w)(T-t)\;t\,dt=\frac1{\beta(T)}\int_0^T\beta'(s)(T-s)\,ds=\frac{1}{\beta(T)}\int_0^T\beta(s)\,ds=\frac{\lambda T+\log(1-\beta(T))}{\mu\,\beta(T)}=\mathbb E[t_{\mathcal P}\mid\bullet\bullet].\]
The second step integrates by parts using \(\beta(0)=0\). For the last,
verify\_all checks \(\frac{d}{dT}[(\lambda T+\log(1-\beta))/\mu]=\beta\)
and that the expression vanishes at \(T=0\). At \(\lambda=\mu\),
\(\int_0^T\beta=T-\mu^{-1}\log(1+\mu T)\), so
\(\mathbb E[t_{\mathcal P}\mid\bullet\bullet]=(1+\mu T)\big(T-\mu^{-1}\log(1+\mu T)\big)/(\mu T)\)
(0.6137 at \(\lambda=\mu=T=1\)).

So the printed (17) has solution
\(Z_{\bullet\bullet}+\mathbb E[t_{\mathcal P}\mid\bullet\bullet]\). This
differs from \(Z_{\bullet\bullet}\) for every \(T>0\). By §2 the
resulting \(\hat Z\) is wrong by
\(N_{\bullet\bullet}\,\mathbb E[t_{\mathcal P}\mid\bullet\bullet]\) for
every alignment with \(N_{\bullet\bullet}\ge1\). The correction is
\[\hat Z_{\bullet\bullet}(T)=\int_0^Tc(t\mid T)\Big(\lambda_v(T-t)\hat Z_\bullet(T-t)+\lambda_w(T-t)\big(\hat Z_\bullet(T-t)+\hat Z_{\circ\bullet}(T-t)\big)\Big)dt ,\]
whose solution is \(Z_\circ\) identically.

\textbf{Evidence} (\texttt{numeric\_check\_output.txt}; checks only):

\begin{enumerate}
\def\labelenumi{\arabic{enumi}.}
\tightlist
\item
  \textbf{Quadrature.} The printed right-hand side was evaluated at 60
  digits at 8 parameter points. It exceeds \(Z_{\bullet\bullet}\) by
  exactly \(\int\beta/\beta\); the residuals are \(\le4\times10^{-18}\),
  and the excesses range from 0.051 to 2.12.
\item
  \textbf{Exact E-step} (\texttt{tkf91\_stats\_from\_counts}, float64,
  200 random paths at each of 4 parameter sets). Corrected (17′) with
  ``\(-T\)'' reproduces \(\mathbb E[S]\) to \(2\times10^{-11}\), and
  \(\mathbb E[B],\mathbb E[D]\) to \(10^{-12}\). Printed (17) overshoots
  by \(N_{\bullet\bullet}\times\int\beta/\beta\), to
  \(2\times10^{-11}\).
\item
  \textbf{{[}DART{]}.} It has
  \texttt{Z11\ +=\ c*(lambdaD*Z1\ +\ lambdaI*(Z1+Z01))}, with no \(t\).
  Its port converges to \(Z_{\bullet\bullet}\) (§4).
\item
  \textbf{Route C Monte Carlo} (whole histories, \(1.5\times10^6\)
  fragments per cell). Ancestral-marker \(\bullet\bullet\) zones under
  the exposure allocation of §5.2 give \(0.4223\pm0.0004\) against
  \(Z_{\bullet\bullet}=0.4224\), at \((\lambda,\mu,T)=(1,0.5,1)\).
  Printed (17) gives 1.018.
\item
  \textbf{Holmes's Fig. 9} (bottom, sequence-length panel; supplement
  p.~14). The ancestor's column is grey (the final \(\bullet\) zone)
  from time 0 to the direct propagation event, cross-hatched (the
  \(\bullet\bullet\) zone) from there to the indirect propagation event,
  and hatched (the \(\circ\bullet\) zone) from there until the ancestor
  dies. So each zone receives the ancestor's exposure only after that
  zone's own propagation event; the exposure before it is in the final
  \(\bullet\) zone. This is the lineage allocation of §5.2, and it
  agrees with (16) and (17′). Printed (17) would also charge the
  \(\bullet\bullet\) zone with the grey part.
\end{enumerate}

\hypertarget{which-accounting-is-consistent-ians-question}{%
\subsection{5.2 Which accounting is consistent (Ian's
question)}\label{which-accounting-is-consistent-ians-question}}

Theorems 1--3 already decide the question without any pathwise
accounting. There is exactly one set of per-zone coefficients for which
Holmes's constant-free sums are exact. It solves (16) as printed and
(17′), not (17).

The pathwise reading is as follows.

\begin{itemize}
\tightlist
\item
  \textbf{Where the marker's exposure goes.} In the generating-function
  representation (§6), \(P(\text{surv},n)\propto e^{-\mu T}W B^n\) and
  \(P(\text{die},n\ge1)\propto W R B^{n-1}\). The factor
  \(W=e^{-\theta T}\mathcal I_D(T,T;0)\) carries the tilt for a full
  \(T\) of carrier exposure, and it sits in the family's final
  \(\bullet\) zone. This is (15)'s ``\(\int_0^T1\,dt\)''.

  \begin{itemize}
  \item
    The marker's exposure before its propagation event is part of that
    \(T\).
  \item
    The \(\circ\bullet\) factor \(R\) and the \(\bullet\bullet\) factor
    \(B\) carry only the marker's exposure after \(t_{\mathcal P}\),
    plus the exposure of extinct insertions.
  \item
    \begin{enumerate}
    \def\labelenumi{(\arabic{enumi})}
    \setcounter{enumi}{15}
    \tightlist
    \item
      follows this accounting; (17) adds the pre-propagation exposure a
      second time.
    \end{enumerate}
  \end{itemize}
\item
  \textbf{Window-union lemma} (route B, proved in prose). Let
  \(X=A_0\to\dots\to A_m=Y\) be the lineage from the marker to the
  survivor \(Y\) nearest the marker on its right. The birth windows
  \((\tau_k,\tau_{k+1})\) of the lineage links, together with
  \((\tau_m,T)\) for \(Y\), partition \([0,T]\). So everything right of
  \(Y\) has the law of the children of a surviving ancestral link over
  \([0,T]\), and the carrier exposure is exactly \(T\). This proves
  ``WLOG the marker is ancestral'' in distribution, under the
  \emph{lineage} allocation: link \(L\)'s exposure at time \(\tau\) is
  charged to the zone in which a child inserted by \(L\) at \(\tau\)
  would sit. Holmes's Fig. 9 (bottom) draws this allocation (§5.1,
  evidence 5), and route C's Monte Carlo confirms it. Under this rule,
  inserted- and immortal-marker zones match ancestral ones:
  \(\max|z|=2.7\) over 144 cells, and 1 of 1710 per-context tests
  exceeds \(|z|=4\).
\item
  \textbf{The alternative, ``naive'', allocation} charges a link's whole
  lifetime to its own zone.

  \begin{itemize}
  \tightlist
  \item
    Under it, the ancestral-marker values would be (16) with a
    ``\(t+\)'' and (17) as printed, with that modified
    \(Z_{\circ\bullet}\) inside.
  \item
    But inserted-marker zones then carry much less exposure: for
    \(\bullet\) at \((1,0.5,1)\), about 0.44--0.47 against 1.047 (0.465
    when the family founder is a mortal ancestral link, 0.440 when it is
    the immortal link). So WLOG fails, and no single set of per-zone
    values makes the sums exact.
  \item
    The printed pair, (16) without \(t\) and (17) with \(t\), fits
    neither rule. Printed (17) is about 50 s.e. from the naive-rule
    Monte Carlo value: 1.0181 against \(1.0289\pm0.0002\).
  \end{itemize}
\end{itemize}

\hypertarget{the-other-items}{%
\subsection{5.3 The other items}\label{the-other-items}}

\begin{itemize}
\tightlist
\item
  \textbf{§5.4 ``\(-T\)''.} See Theorem 1 and the uniqueness paragraph
  of §2. Numerically, the zone sum minus \(\mathbb E[S]\) is \(T\) to
  \(2\times10^{-11}\) on every path
  (\texttt{numeric\_check\_output.txt}, part 2).
\item
  \textbf{Fig. 9 caption.} In the lower drawing the ancestor's line ends
  in a cross, then a dotted segment, then an open circle at \(T\), and
  the caption itself lists a \(\circ\bullet\) zone first. So the
  ancestor dies, and ``Ancestor survives'' is wrong.
\item
  \textbf{\(q_0\).} See A1. The text before (11) also says
  \(\alpha(t)q_0(T-t)\) is ``not deleted during \([0,t]\) and becomes
  extinct over \([t,T]\)'', which fixes it.
\item
  \textbf{{[}DART{]} implements the corrected system.} In
  \texttt{TKF\_coeff::initialise}, line 131 updates \texttt{Z01} by
  (16), and line 132 updates \texttt{Z11} by
  \(c\,(\lambda_vZ_\bullet+\lambda_w(Z_\bullet+Z_{\circ\bullet}))\),
  i.e.~(17′) with no ``\(t+\)'' (\texttt{lambdaD}, \texttt{lambdaI},
  \texttt{lambdaE} are \(\lambda_v,\lambda_w,\lambda_x\)). In
  \texttt{accumulate\_counts}, line 178 is \texttt{deltaZ\ -=\ T}
  (``subtract site-time contribution of immortal link''), the ``\(-T\)''
  of misprint 2. \texttt{set\_params} (line 205) sets
  \(\lambda=I/(Z+T)\), \(\mu=D/Z\).
\item
  \textbf{Checked and correct:}

  \begin{itemize}
  \item
    \begin{enumerate}
    \def\labelenumi{(\arabic{enumi})}
    \setcounter{enumi}{7}
    \tightlist
    \item
      and (13);
    \end{enumerate}
  \item
    (9)--(16);
  \item
    \(\hat S=N_\circ+N_\bullet-1\), which is the ancestral length \(i\)
    (Lemma 1);
  \item
    \(\hat D=\hat I+N_\circ-N_{\bullet\bullet}\);
  \item
    the zone-type lists of §5.4.
  \end{itemize}
\end{itemize}

\hypertarget{relation-to-eqvolterra-q2}{%
\section{6. Relation to eq:volterra
(Q2)}\label{relation-to-eqvolterra-q2}}

\hypertarget{tilted-grammar}{%
\subsection{6.1 Tilted grammar}\label{tilted-grammar}}

Weight each parse tree by \(e^{-\theta S_{\rm fam}}\), where
\(S_{\rm fam}\) includes the founding link's exposure. Only the \(L\)
productions change:
\[L(t)\to A\,D(t,t):\ e^{-(\mu+\theta)t},\qquad L(t)\to G_\ell D(t,\ell):\ \mu e^{-(\mu+\theta)\ell}d\ell .\]
Write
\(F_t=F^{\theta}_t(z,w)=\mathbb E[z^{N_a}w^{N_g}e^{-\theta S_{\rm fam}}]\).
The derivation of app:pgf goes through unchanged:
\[\mathcal I^\theta_D(t,T)=\exp\Big(\lambda\!\int_{t-T}^t(F_s-1)ds\Big),\qquad
F_t=z\,e^{-(\mu+\theta)t}\mathcal I^\theta_D(t,t)+w\!\int_0^t\!\mu e^{-(\mu+\theta)\ell}\,\mathcal I^\theta_D(t,\ell)\,d\ell,\tag{\ensuremath{\mathrm V_\theta}}\]
\[\partial_tF_t=\lambda F_t^2-(\lambda+\mu+\theta)F_t+\mu w,\qquad F_0=z.\tag{\ensuremath{\mathrm R_\theta}}\]

This is eq:volterra and eq:riccati with
\((\mu,w)\mapsto(\mu+\theta,\mu w/(\mu+\theta))\). So eq:pgf-closed and
eq:ID-closed hold with \(z_\pm\) the roots of
\(\lambda x^2-(\lambda+\mu+\theta)x+\mu w\), and
\(\Delta=\lambda(z_+-z_-)\).

eq:pgf-closed takes \(z_+\) to be the larger root; we do not. All closed
forms are invariant under \(z_+\leftrightarrow z_-\) with
\(\Delta\to-\Delta\) (then \(\rho_s\to1/\rho_s\)). At
\((\theta,w)=(0,1)\) take \(z_+=\mu/\lambda\), \(z_-=1\); this is valid
for either sign of \(\mu-\lambda\). verify\_all §5 parametrises by the
roots, so no square root appears.

\hypertarget{lemmas}{%
\subsection{6.2 Lemmas}\label{lemmas}}

Let \(E=e^{\Delta T}\). All statements hold at general \((\theta,w)\);
verify\_all §5 checks each of (a)--(g) there.

\begin{enumerate}
\def\labelenumi{(\alph{enumi})}
\item
  \(F_T(0)=\frac{\mu w}{\lambda}B\), where \(B:=\frac{E-1}{z_+E-z_-}\).
  Both sides satisfy \(y'=\mu w-(\lambda+\mu+\theta)y+\lambda y^2\) with
  \(y(0)=0\).
\item
  \(\mathcal I_D(T,T;z)=\mathcal I_D(T,T;0)/(1-Bz)\). Hence
  \([z^n]\,e^{-\theta T}\mathcal I_D(T,T;z)=W B^n\), with
  \(W:=e^{-\theta T}\mathcal I_D(T,T;0)=e^{(\lambda z_+-\lambda-\theta)T}(z_+-z_-)/(z_+E-z_-)\).
\item
  \(F^{(1)}:=[z^1]F_T=(1-z_+B)(1-z_-B)\), and
  \(\partial_TB=\lambda F^{(1)}\) with \(B(0)=0\).
\item
  \textbf{Additivity.}
  \(\mathcal I_D(T,t;0)\,e^{-\theta t}\,W(T-t)=W(T)\). \emph{Proof:}
  \(\mathcal I_D(T,t;0)=\exp(\lambda\int_{T-t}^T(F_s(0)-1)ds)\) and
  \(W(T-t)=e^{-\theta(T-t)}\exp(\lambda\int_0^{T-t}(F_s(0)-1)ds)\); the
  product with \(e^{-\theta t}\) joins the two ranges into \([0,T]\).
\item
  \textbf{Survive and die parts.} \([z^1]\) of the \(A\)-term is
  \(e^{-\mu T}W\). Put \(R:=[z^1]G_T/W\), where \(G_T\) is the
  \(G_\ell\)-term. Then \(F^{(1)}=W(e^{-\mu T}+R)\), and
  \[R=\frac{(z_+-z_-)\,e^{(\lambda+\theta-\lambda z_-)T}}{z_+E-z_-}-e^{-\mu T}.\]
\item
  \textbf{Values at \((0,1)\).} \(B=\beta\), \(W=1-\beta\),
  \(R=(1-\alpha)\gamma\) and \(F_T(0)=q\).
\item
  \textbf{Semigroup (S).}
  \(w\int_t^T\mu e^{-(\mu+\theta)\ell}\mathcal I_D(T,\ell;0)d\ell=e^{-(\mu+\theta)t}\mathcal I_D(T,t;0)F_{T-t}(0)\).
  At \((0,1)\) this reads
  \(\int_t^T\mu v(\ell\mid T)d\ell=v(t\mid T)q(s)\). \emph{Proof:} both
  sides vanish at \(t=T\), since \(F_0(0)=0\). On the right,
  \(\partial_t\log\mathcal I_D(T,t;0)=\lambda(F_{T-t}(0)-1)\), and
  (R\(_\theta\)) gives
  \(\partial_tF_{T-t}(0)=-(\lambda F^2-(\lambda+\mu+\theta)F+\mu w)\) at
  \(F=F_{T-t}(0)\). The terms in \(F\) and \(F^2\) cancel, leaving
  \(-\mu w\,e^{-(\mu+\theta)t}\mathcal I_D(T,t;0)\), which is the
  \(t\)-derivative of the left side.
\end{enumerate}

\textbf{Proposition 6.} At \((\theta,w)=(0,1)\): \[\begin{aligned}
I_\circ&=\partial_w\log B, & Z_\circ&=-\partial_\theta\log B, &
I_\bullet&=\partial_w\log W, & Z_\bullet&=-\partial_\theta\log W,\\
I_{\circ\bullet}&=\partial_w\log R, & Z_{\circ\bullet}&=-\partial_\theta\log R, &
I_{\bullet\bullet}&=1+\partial_w\log B, & Z_{\bullet\bullet}&=-\partial_\theta\log B .
\end{aligned}\]

These are verified symbolically against Definition 2, using the implicit
derivatives of \(z_\pm\) (verify\_all §5, ``GF route == score route'').
Those derivatives are singular at \(\lambda=\mu\), \((\theta,w)=(0,1)\),
where the roots coincide, so the verification covers \(\lambda\ne\mu\).
Proposition 6 extends to \(\lambda=\mu\) by continuity: \(B\), \(W\) and
\(R\) are symmetric in \(z_\pm\), hence analytic in \(\Delta^2\), and
Definition 2 extends by its \(k\to0\) limits (§2). The probabilistic
reading:

\begin{itemize}
\tightlist
\item
  \(P(\text{die},0)=F_T(0)=(\mu w/\lambda)B\),
  \(P(\text{surv},n)=e^{-\mu T}WB^n\) and \(P(\text{die},n)=WRB^{n-1}\).
\item
  The family table of §2 then gives the formulas, since
  \(\partial_w\log(\mu w/\lambda)=1\) and
  \(\partial_\theta\log(\mu w/\lambda)=0\).
\item
  This also explains \(I_{\bullet\bullet}=I_\circ+1\) and
  \(Z_{\bullet\bullet}=Z_\circ\) (§3): the extinction PGF and the
  geometric ratio of the survivor count differ by the \(\theta\)-free
  factor \(\mu w/\lambda\), whose \(\partial_w\log\) is 1.
\end{itemize}

\hypertarget{holmess-equations-as-variations-of-v_theta}{%
\subsection{\texorpdfstring{6.3 Holmes's equations as variations of
(V\(_\theta\))}{6.3 Holmes's equations as variations of (V\_\textbackslash{}theta)}}\label{holmess-equations-as-variations-of-v_theta}}

\textbf{(9) and (14): linearisation at \(z=0\).}

\begin{itemize}
\tightlist
\item
  At \(z=0\), (V\(_\theta\)) is
  \(F_T(0)=w\int_0^T\mu e^{-(\mu+\theta)\ell}\exp\big(\lambda\int_{T-\ell}^T(F_s(0)-1)ds\big)d\ell\).
\item
  Apply \(\partial_w\) at \((0,1)\) and use
  \(\mu e^{-\mu\ell}\mathcal I_D(T,\ell;0,1)=\mu v(\ell\mid T)\):
  \[\partial_wF_T(0)=q(T)+\int_0^T\mu v(\ell\mid T)\,\lambda\!\int_{T-\ell}^T\partial_wF_{s}(0)\,ds\,d\ell .\]
\item
  Exchange the order of integration on \(\{0<t<\ell<T\}\), with
  \(s=T-t\). Apply (S), substitute
  \(\partial_wF_s(0)=q(s)(1+I_\circ(s))\) (Proposition 6), and divide by
  \(q(T)\). This gives exactly (9), with \(a=vq(s)/q(T)\) and
  \(\lambda_x=\lambda q\).
\item
  \(-\partial_\theta\) gives (14) in the same way. The factor
  \(e^{-\theta\ell}\) contributes \(\ell=\int_0^\ell dt\), which becomes
  the ``1'' in (14).
\end{itemize}

The variation of a nonlinear Volterra equation \(F=\mathcal N[F]\) is
the linear equation
\(\delta F=\partial\mathcal N+D\mathcal N[F]\,\delta F\), a Volterra
equation of the second kind. Its kernel is
\(D\mathcal N[F(\cdot,0)]\,h\,(T)=\int_0^T\lambda v(t\mid T)q(s)h(s)\,dt\).
Written for \(h=q\,u\), this is
\(q(T)\int_0^Ta(t\mid T)\lambda_x(s)u(s)\,dt\), which is Holmes's
kernel.

\textbf{(11) and (16): the \(z^1\) coefficient.}

\begin{itemize}
\tightlist
\item
  \([z^1]\) of the \(G_\ell\) term, after Fubini and (S), gives
  \[[z^1]G_T=\int_0^T\lambda F^{(1)}_{s}\;e^{-(\mu+\theta)t}\mathcal I_D(T,t;0)\,F_{s}(0)\,dt .\]
  The factors read:

  \begin{itemize}
  \tightlist
  \item
    \(t\) is the time of the most recent birth by the marker whose
    family contains a survivor (in the \(z^1\) coefficient, the only
    survivor, which is the survivor nearest the marker on its right).
    This is Holmes's propagation event, ``the last (directly or
    indirectly) surviving insertion from the ancestral link'';
  \item
    \(\mathcal I_D(T,t;0)\): earlier births, all extinct;
  \item
    \(F_s(0)\): the marker after \(t\), dead with no survivor.
  \end{itemize}
\item
  Divide by \(W(T)\) using (d), and use (e):
  \[R(T)=\int_0^T\lambda\,e^{-\mu t}\,\big[e^{-\mu s}+R(s)\big]\,F_s(0)\,dt .\tag{\ensuremath{\star}}\]
  verify\_all checks (\(\star\)) at general \((\theta,w)\) in the
  differential form
  \(\frac{d}{dT}[Re^{\mu T}]=\lambda(1+e^{\mu T}R)F_T(0)\), \(R(0)=0\).
\item
  The factor \(e^{-\mu t}\) is untilted. The marker's tilt
  \(e^{-\theta t}\) before \(t_{\mathcal P}\) has moved into \(W(T)\),
  i.e.~into the final \(\bullet\) zone.
\item
  Apply \(\partial_w\) to (\(\star\)), and divide by
  \(R(T)=(1-\alpha)\gamma\) (so
  \(\lambda\alpha(t)q(s)/R(T)=\lambda b(t\mid T)\), and
  \(e^{-\mu s}+R(s)=1-q(s)\) at \((0,1)\)). This gives (11). Applying
  \(-\partial_\theta\) instead gives (16), with no \(t\).
\end{itemize}

\textbf{(10), (15), (12) and (17′): quadratures.}

\begin{itemize}
\tightlist
\item
  \(W=e^{-\theta T}\exp(\lambda\int_0^T(F_s(0)-1)ds)\), so
  \(\partial_w\log W\) and \(-\partial_\theta\log W\) are (10) and (15).
\item
  \(B=\lambda\int_0^TW(s)(e^{-\mu s}+R(s))ds\), by (c) and (e). At
  \((0,1)\), \(W(s)/\beta(T)=c(t\mid T)\) and
  \(\int_0^Tc\,\lambda(1-q(s))dt=1\). Then \(1+\partial_w\log B\) is
  (12), and \(-\partial_\theta\log B\) is (17′).
\item
  No term in \(B\) carries the marker's pre-propagation tilt. Inserting
  one reproduces the printed (17).
\end{itemize}

\hypertarget{propagation-event-and-dto-lcdot-d}{%
\subsection{\texorpdfstring{6.4 Propagation event and
\(D\to L\cdot D\)}{6.4 Propagation event and D\textbackslash{}to L\textbackslash{}cdot D}}\label{propagation-event-and-dto-lcdot-d}}

\textbf{The grammar's production.} eq:gw-D,
\(D(t,T)\to L(t-u)\cdot D(t,u)\) with density
\(\lambda e^{-\lambda(T-u)}\), emits the \emph{most recent} birth in the
window first; it is placed leftmost. The remaining window \((0,u)\) is
left to \(D(t,u)\).

\textbf{Holmes's propagation event} is the most recent birth whose
family has a survivor. Split each \(L\) by \(z\)-degree. The children
born after \(\mathcal P\) are all extinct (\(z^0\), weight \(F(0)=q\)).
Summing the \(D\to L\cdot D\) recursion over them replaces
\(\lambda e^{-\lambda(T-u)}\) by
\(\lambda e^{-\int\lambda(1-q)}=\lambda\mathcal I_D(\cdot;0)\). Over the
window after \(\mathcal P\) this gives
\(\mathcal I_D(T-t,T-t;0)=1-\beta(T-t)\), the numerator of
\(c(t\mid T)\). The same functional form over the earlier window
\([0,t]\) is \(v(t\mid T)/e^{-\mu t}=\mathcal I_D(T,t;0)\) in (8); it
enters \(a(t\mid T)\), not \(b\).

Analytically, this is the first-passage identity
\[e^{\lambda\int_0^T\tilde F}-1=\int_0^T\lambda\tilde F_{T-t}\,e^{\lambda\int_0^t\tilde F_{T-u}du}\,dt,\qquad \tilde F=F-F(0),\]
applied inside \(\mathcal I_D\). It gives \(B=\lambda\int_0^TF^{(1)}\),
i.e.~(c). So the two decompositions coincide once the extinct children
are grouped.

The split into direct and indirect propagation, \(\lambda_v\) and
\(\lambda_w\), is the split of \(F^{(1)}=W(e^{-\mu s}+R)\) into its
\(A\) and \(G_\ell\) parts.

\textbf{(9) and (14) differ.} They do not decompose on the last event.
They take the first moment, over all extinct births, of the Poisson
cloud generated by \(\mathcal I_D\). They also exchange the grammar's
order (death time \(\ell\) first, in \(G_\ell\)) with Holmes's order
(birth time \(t\) first). In the factor
\(\int_t^T\mu v\,d\ell=v\,q(s)\),
\(q(s)=(1-\alpha(T-t))(1-\gamma(T-t))\) is the factor on the marker
after \(t\) in Holmes's Fig. 8(i).

\hypertarget{where-the-three-routes-disagreed}{%
\section{7. Where the three routes
disagreed}\label{where-the-three-routes-disagreed}}

There were no substantive disagreements. All three routes conclude the
following:

\begin{itemize}
\item
  (8)--(16) and (13) are correct;
\item
  \begin{enumerate}
  \def\labelenumi{(\arabic{enumi})}
  \setcounter{enumi}{16}
  \tightlist
  \item
    loses its ``\(t+\)'' terms, with the same excess;
  \end{enumerate}
\item
  Holmes's zone sum equals \(\mathbb E[S]+T\);
\item
  \(q_0=(1-\alpha)(1-\gamma)\).
\end{itemize}

These conclusions are settled here algebraically (§§2--5) and against
{[}DART{]}.

The differences in detail, and how they are resolved:

\begin{enumerate}
\def\labelenumi{\arabic{enumi}.}
\tightlist
\item
  \textbf{The ``\(+T\)''.} Routes A and B concluded that ``Holmes's
  \(Z\) includes the immortal link'' and left his main text unchecked.
  {[}H-M{]} §3.3 defines \(\hat Z\) as mortal-link time, and {[}DART{]}
  subtracts \(T\). So the supplement's §5.4 sum is short of a ``\(-T\)''
  (misprint 2). Holmes's M-step is \(\lambda=\hat I/(\hat Z+T)\),
  \(\mu=\hat D/\hat Z\), which is eq:mstep.
\item
  \textbf{Fig. 9 caption.} Only route B checked it. I confirmed it from
  the page image.
\item
  \textbf{Zone weight conventions.} Route A uses \(R=(1-\alpha)\gamma\)
  for \(\circ\bullet\) with the marker tilt absorbed into \(W\). Route B
  uses \(P_{\circ\bullet}=(1-\alpha)\gamma/\alpha\) plus the constant
  \(\alpha^{-1}\). Both give the same \(I_z,Z_z\) (verify\_all §5,
  Proposition 6), so this is bookkeeping only.
\item
  \textbf{Routes A, B and C state the closed forms in different
  variables} (\(E=e^{(\mu-\lambda)T}\); \((\eta,\alpha)\); \((x,y)\)).
  All are checked equal to Definition 2.
\item
  \textbf{Uniqueness.} Route A cited Linz; route B used Gronwall for (9)
  and a contraction argument for (11). Theorem 3 reduces all four
  homogeneous equations to \(u(T)=\int_0^Tk\,u\) with
  \(0\le k\le\lambda\), so one Gronwall step covers (11) and (16) as
  well, with no contraction argument.
\item
  \textbf{Accounting.} Route C named the exposure allocations ``naive''
  and ``lineage''; route B's window-union lemma is the lineage rule.
  Neither is needed for Theorems 1--3.
\item
  \textbf{Status of route C's ``proved as printed''.} For (8) and (13)
  this was numerical only. Here they are proved symbolically.
\end{enumerate}

\hypertarget{open-issues}{%
\section{8. Open issues}\label{open-issues}}

\begin{itemize}
\tightlist
\item
  \textbf{\(q_0\).} A1 is not restated in the supplement; it is Eq. (2)
  of {[}H-M{]} at \(n=0\), and agrees with (13), Fig. 8 and the prose.
\item
  \textbf{The window-union lemma (§5.2) is proved in prose only.} It
  explains the pathwise accounting but is not needed for Theorems 1--3.
  Its consequences are confirmed by route C's Monte Carlo.
\item
  \textbf{Series dependence.} Initial conditions and the
  \(\lambda\to\mu\) limits rely on sympy series expansions. These are
  symbolic computations, not floating point.
\item
  \textbf{Smoothness in \((\theta,w)\).} In §6.3, differentiating under
  the integral in \((\theta,w)\) uses smoothness of the closed forms
  near \((0,1)\) for finite \(T\); this was not written out. §6 is in
  any case an explanation (Q2). Q1 and Q3 rest on §§2--4, which do not
  use it.
\item
  \textbf{The \(\lambda\leftrightarrow\mu\) symmetry} of all zone
  coefficients is verified but not explained.
\item
  \textbf{{[}H-M{]} §3.3 formulas.} Its displayed M-step formulas are
  images in the PMC copy. Misprint 2 rests on its prose definition of
  \(\hat Z\) (``expected total time accrued by mortal links; \(T\)
  \ldots{} by the immortal link'') and on {[}DART{]}.
\end{itemize}

\hypertarget{files}{%
\section{Files}\label{files}}

\begin{itemize}
\tightlist
\item
  \texttt{proof.md}, \texttt{proof.tex}: this note.
\item
  \texttt{verify\_all.py} → \texttt{verify\_all\_output.txt}: 172
  symbolic checks, all passing; runs in about 60 s.
\item
  \texttt{numeric\_check.py} → \texttt{numeric\_check\_output.txt}:
  three numerical checks for Q3 --- quadrature of printed (17), the
  exact E-step (\texttt{tkf91\_stats\_from\_counts}), and the {[}DART{]}
  port.
\end{itemize}

\end{document}
